\documentclass[letter, 10pt]{article}
\usepackage{float}
\usepackage[utf8]{inputenc}
\usepackage[spanish]{babel}
\usepackage{amsfonts}
\usepackage{amsmath}
\usepackage[dvips]{graphicx}
\usepackage{url}
\usepackage[top=3cm,bottom=3cm,left=3.5cm,right=3.5cm,footskip=1.5cm,headheight=1.5cm,headsep=.5cm,textheight=3cm]{geometry}



\begin{document}
\title{Inteligencia Artificial \\ \begin{Large}Estado del Arte: Problema [Nombre Problema]\end{Large}}
\author{[Nombre autor]}
\date{\today}
\maketitle


%--------------------No borrar esta secci\'on--------------------------------%
\section*{Evaluaci\'on}

\begin{tabular}{ll}
Resumen (5\%): & \underline{\hspace{2cm}} \\
Introducci\'on (5\%):  & \underline{\hspace{2cm}} \\
Definici\'on del Problema (10\%):  & \underline{\hspace{2cm}} \\
Estado del Arte (40\%):  & \underline{\hspace{2cm}} \\
Modelo Matem\'atico (20\%): &  \underline{\hspace{2cm}}\\
Conclusiones (15\%): &  \underline{\hspace{2cm}}\\
Bibliograf\'ia (5\%): & \underline{\hspace{2cm}}\\
 &  \\
\textbf{Nota Final (100\%)}:   & \underline{\hspace{2cm}}
\end{tabular}
%---------------------------------------------------------------------------%
\vspace{2cm}

\begin{abstract}
Resumen del informe en no m\'as de 10 l\'ineas, donde se sintetice el problema que se trata y sirva para que un lector no involucrado comprenda el objetivo del documento.

\end{abstract}

\section{Introducci\'on}
%Una explicaci\'on breve del contenido del informe, es decir, detalla: Prop\'osito, Estructura del Documento, Descripci\'on (muy breve) del Problema y Motivaci\'on.

En este informe se realizará una presentación de los trabajos más importantes respecto al problema de Truck Drone Parcel Delivery Problem (TDPDP), el cual consiste en encontrar una ruta para un camión que haga entregas a ciertas ubicaciones en el menor tiempo posible, teniendo el apoyo de un dron que permite hacer cargas en paralelo o con menor tiempo respecto al camión.
La motivación de este trabajo nace de análisis como \cite{moshrefjavadi2021applications} donde se explica el gran auge que ha tenido la utilización de drones en el último tiempo, especialmente con motivos de entrega de productos, dado su alto valor económico para las empresas.
A continuación, se presentarán los principales avances desde su concepción inicial en el trabajo de Murray et al. (2014) \cite{Murray2015Flying}, pasando también por variantes definidas en otros trabajos, y comparando de qué manera se resolvían estos problemas en computación.


\section{Definici\'on del Problema}
%Explicaci\'on del problema que se va a estudiar, en qu\'e consiste, cu\'ales son sus variables , restricciones y objetivo(s) de manera general (en palabras, no una formulaci\'on matem\'atica). Debe entenderse claramente el problema y qu\'e busca resolver.
%Explicar si existen problemas relacionados.
%Destacar, si existen, las variantes m\'as conocidas.\\
%Redactar en tercera persona, sin faltas de ortograf\'ia y referenciar correctamente sus fuentes mediante el comando  \verb+\cite{ }+. Por ejemplo, para hacer referencia al art\'iculo de algoritmos h\'ibridos para problemas de satisfacci\'on
%de restricciones~\cite{Prosser93Hybrid}.

El problema de Truck Drone Parcel Delivery Problem (Que nos referiremos a el de la forma TDPDP) ha tenido varias definiciones en varios trabajos. Antes de la definición que usaremos durante este trabajo, haremos referencias a otras definiciones realizadas en trabajos previos.

Primero, hemos de mencionar que TDPDP es una variante de otros problemas más conocidos en el mundo de la optimización. Estos corresponden al TSP (The Salesman Problem, o El problema del Vendedor Viajero en español) y el problema de VRP (Vehicle Routing Problem, o el Problema de Ruteo de Vehiculos en español). VRP es la variante más cercana a TDPDP, y la cual es también variante de TSP, por lo que nos enfocaremos en referirnos a esta como el problema de donde proviene TDPDP.

VRP , definido según \cite{Irnich2014FamilyVRP}, se tiene que \textbf{dado} un set de entregadas solicitadas (también entendidas como ubicaciones) y una cantidad de vehiculos, entonces la \textbf{tarea} es determinar un set de ruta para cada uno de los vehiculos para poder llegar a las ubicaciones de entrega a un costo mínimo, es decir, decidir que vehiculo va a cuales ubicaciones y en que secuencia para que todas (o algunas) de las ubicaciones sean cubiertas. 

De la definición de VRP, se define una variante conocida como CVRP, la cual se acerca más a nuestra definición para el TDPDP. Esta variante tiene un punto de inicio (conocido como depot, deposito en español), y se le hace también referencia a puntos (o ubicaciones) de entrega. El problema debe resolverse desde el punto de inicio definido hacia cada una de las ubicaciones de entrega. Las restricciones del problema siguen a VRP, es decir, se tienen también la restricción de visitar todas las ubicaciones de entrega.

Conociendo el problema de VRP, TDPDP es una variante de VRP en la que se incluye un dron, el cual permite realizar entregas en conjunto al vehiculo. Siguiendo la definición realizada en \cite{Murray2015Flying}, este dron tiene las caracteristicas de que su costo para viajar de un nodo puede ser distinto que al costo del camión y, adicionalmente, hacer entregas en paralelo mientras el camión visita otros nodos, lo cual puede hacer el costo asociado al viaje más bajo y los tiempos de entrega más cortos. El objetivo es minimizar el tiempo de retorno más tardio de ambos vehiculos al deposito.

Se tienen restricciones relacionadas con el dron. Según \cite{Murray2015Flying}, hay restricciones con respecto a la cantidad de drones por camión, la bateria del dron, y las velocidades del dron. Solo puede haber un dron asociado a cada camión, el cual al salir del camión debe de reunirse con este en algún nodo posterior. El dron tiene una bateria asociada la cual indica cuanta distancia puede recorrer y cuanto tiempo puede quedarse en espera (Por ejemplo, un dron podria esperar en un nodo al camión siempre y cuando le quede bateria cuando este llegue), y por ultimo la velocidad del dron limita que nodos pueden visitarse.




\section{Estado del Arte}
La informaci\'on que describen en este punto se basa en los estudios realizados con antelaci\'on respecto al tema.
Lo m\'as importante que se ha hecho hasta ahora con relaci\'on al problema. Deber\'ia responder preguntas como las siguientes:
?`cu\'ando surge?, ?`qu\'e m\'etodos se han usado para resolverlo?, ?`cu\'ales son los mejores algoritmos que se han creado hasta
la fecha?, ?`qu\'e representaciones han tenido los mejores resultados?, ?`cu\'al es la tendencia actual para resolver el problema?, tipos de movimientos, heur\'isticas, m\'etodos completos, tendencias, etc... Puede incluir gr\'aficos comparativos o explicativos.\\



\section{Modelo Matem\'atico}
%Uno o m\'as modelos matem\'aticos para el problema, idealmente indicando el espacio de b\'usqueda para cada uno. Cada modelo debe estar correctamente referenciado, adem\'as no debe ser una imagen extraida. Tambi\'en deben explicarse en detalle cada una de las partes, mostrando claramente la funci\'on a maximizar/minimizar, variables y restricciones. Tanto las f\'ormulas como las explicaciones deben ser consistentes.

Para definir el modelo matemático utilizaremos varias de las fuentes. Los simbolos utilizados referenciaran directamente el trabajo de \cite{Irnich2014FamilyVRP}, pero hablaremos de las diferencias con los otros modelos en su sección pertinente.

\subsection{Parametros}

Se definen los siguientes parametros:
\begin{itemize}
\item $C \in \{1,\dots,c\}$: Conjunto de clientes
\item $C' \subset C$: Conjunto de clientes que pueden ser entregados por el dron.
\item $N = \{0,1,\dots,c+1\}$ Conjunto de todos los nodos (o ubicaciones de entrega) posibles. $0$ indica el depot inicial y $c+1$ volver al depot.
\item $N_0 \in \{0,1,\dots,c\}$: Conjunto de todos los nodos donde el vehiculo puede partir.
\item $N_+ \in \{1,\dots,c,c+1\}$: Conjunto de todos los nodos donde el vehiculo puede visitar.
\item $\tau_{ij}$: Tiempo requerido por el camión para viajar del nodo $i$ al nodo $j$
\item $\tau'_{ij}$: Tiempo requerido por el dron para viajar del nodo $i$ al nodo $j$.
\item $S_L$: Tiempo que tarda en lanzarse el dron de alguna ubicación.
\item $S_R$: Tiempo que tarda en recuperarse el dron de alguna ubicación.
\item $e$: Resistencia de vuelo del dron. Indica cuanto tiempo puede soportar volando.
\item $P$: Conjunto de tuplas en la forma $i,j,k$ en donde se incluye si cumple las siguientes condiciones:
\begin{enumerate}
\item El punto de salida $i$ no debe ser un nodo de termino ($i \in N_0$)
\item El punto de llegada $j$ debe de no ser igual al punto $i$ y debe de pertencer a los nodos que pueden ser entregados por dron. (Es decir, $j \in \{C': j \neq i\}$).
\item El punto de recuperación del dron, $k$, no debe de ser ni $i$ ni $j$, puede ser cualquier nodo de cliente o el depot, y el tiempo de vuelo de $i\to j\to k$ no debe de ser más que la capacidad de resistencia del dron $e$.
\begin{itemize}
\item Matematicamente, se define como $k\in\{N_+: k \neq j, k\neq i, \tau'_{ij}+\tau'_{jk}\leq e\}$
\end{itemize}
\end{enumerate}
\end{itemize}

\subsection{Variables}

Se definen las siguientes variables de decisión:

\begin{itemize}
\item $x_{ij}$: Indica si el camión se mueve del nodo $i$ al nodo $j$.
\item $y_{ijk}$: Indica si el dron sale del nodo $i$, atiende al nodo $j$ y se queda en $k$
\item $t_j$: Indica el tiempo de entrega del camión en el nodo $j$.
\item $t'_j$: Indica el tiempo de entrega del dron en el nodo $j$.
\end{itemize}

El trabajo de \cite{Irnich2014FamilyVRP} también hace referencia al uso de dos variables auxiliares
\begin{itemize}
\item $p_{ij}$: Vale uno si indica que el cliente $i \in C$ es visitado en algún tiempo previo a $j\in\{C: j\neq i\}$.
\item $u_i$: Especifica la position del nodo $i \in N_+$ en el camino del camión. Usada para evitar subrutas del camión, y asi hacer de que si un dron va hacia un nodo el camión no este obligado a pasar por ahi.
\end{itemize}

En donde los indices se mueven en:
\begin{itemize}
\item $i \in N_0$: Nodos de donde se puede partir.
\item $j \in N_+$: Nodos donde se pueden viajar.
\item $k \in P$: Nodos viajables por el dron
\end{itemize}

Los dominios de la variable son:

\begin{itemize}
\item $D_{x_{ij}} \in \{0,1\}$
\item $D_{y_{ijk}} \in \{0,1\}$
\item $D_{t_j} \in \mathbf{N}$
\item $D_{t'_j} \in \mathbf{N}$
\end{itemize}

\subsection{Restricciones}

\subsection{Variaciones}

Como hemos dicho anteriormente, toda la sección anterior hacen directa referencia al \cite{Irnich2014FamilyVRP}. Usaamos este modelo dado de que comparte similitud con la mayoria de los modelos. En esta sección, hablaraemos de algunas variaciones mencionadas en otros modelos.

Cada cliente debe ser visitado exactamente una vez, ya sea por el camión o por el dron:
\\[\sum_{i \in N_0,\, i \neq j} x_{ij} + \sum_{i \in N_0,\, i \neq j} \sum_{\substack{k \in N^+ \\ \langle i,j,k \rangle \in P}} y_{ijk} = 1 \quad \forall j \in C \tag{2}\\]

El camión debe salir una sola vez del depósito de origen y llegar una sola vez al depósito final:
\\[\sum_{j \in N^+} x_{0j} = 1 \tag{3}\\]
\\[\sum_{i \in N_0} x_{i, c+1} = 1 \tag{4}\\]

Conservación de flujo en cada parada de cliente:
\\[\sum_{i \in N_0,\, i \neq j} x_{ij} = \sum_{k \in N^+,\, k \neq j} x_{jk} \quad \forall j \in C \tag{6}\\]

Eliminación de subtours para el recorrido del camión (formulación tipo Miller-Tucker-Zemlin):
\\[u_i - u_j + 1 \le (c+2)(1 - x_{ij}) \quad \forall i \in C, \; j \in N^+ \setminus \{i\} \tag{5}\\]

A lo más un lanzamiento y a lo más una recuperación del dron en cada nodo:
\\[\sum_{j \in C,\, j \neq i} \sum_{\substack{k \in N^+ \\ \langle i,j,k \rangle \in P}} y_{ijk} \le 1 \quad \forall i \in N_0 \tag{7}\\]
\\[\sum_{i \in N_0,\, i \neq k} \sum_{\substack{j \in C \\ \langle i,j,k \rangle \in P}} y_{ijk} \le 1 \quad \forall k \in N^+ \tag{8}\\]

Si el dron realiza una salida \\(\langle i, j, k \rangle\\), el camión debe visitar obligatoriamente los nodos \\(i\\) y \\(k\\):
\\[2 y_{ijk} \le \sum_{h \in N_0,\, h \neq i} x_{hi} + \sum_{l \in C,\, l \neq k} x_{lk} \quad \forall i \in C, \; j \in C \setminus \{i\}, \; k \in \{N^+ : \langle i,j,k \rangle \in P\} \tag{9}\\]

Caso particular cuando el dron es lanzado directamente desde el depósito inicial (\\(i=0\\)):
\\[y_{0jk} \le \sum_{h \in N_0,\, h \neq k} x_{hk} \quad \forall j \in C, \; k \in \{N^+ : \langle 0,j,k \rangle \in P\} \tag{10}\\]

El nodo de lanzamiento \\(i\\) debe preceder al nodo de recuperación \\(k\\) en el orden de paradas del camión:
\\[u_k - u_i \ge 1 - (c+2)\left(1 - \sum_{\substack{j \in C \\ \langle i,j,k \rangle \in P}} y_{ijk}\right) \quad \forall i \in C, \; k \in N^+ \setminus \{i\} \tag{11}\\]

Coordinación temporal en el despegue (fuerza \\(t'_i = t_i\\) en el cliente de lanzamiento \\(i\\)):
\\[t'_i \ge t_i - M\left(1 - \sum_{j \in C,\, j \neq i} \sum_{\substack{k \in N^+ \\ \langle i,j,k \rangle \in P}} y_{ijk}\right) \quad \forall i \in C \tag{12}\\]
\\[t'_i \le t_i + M\left(1 - \sum_{j \in C,\, j \neq i} \sum_{\substack{k \in N^+ \\ \langle i,j,k \rangle \in P}} y_{ijk}\right) \quad \forall i \in C \tag{13}\\]

Coordinación temporal en el reencuentro (fuerza \\(t'_k = t_k\\) en el nodo de llegada \\(k\\)):
\\[t'_k \ge t_k - M\left(1 - \sum_{i \in N_0,\, i \neq k} \sum_{\substack{j \in C \\ \langle i,j,k \rangle \in P}} y_{ijk}\right) \quad \forall k \in N^+ \tag{14}\\]
\\[t'_k \le t_k + M\left(1 - \sum_{i \in N_0,\, i \neq k} \sum_{\substack{j \in C \\ \langle i,j,k \rangle \in P}} y_{ijk}\right) \quad \forall k \in N^+ \tag{15}\\]

Tiempo de llegada del camión al nodo \\(k\\) considerando el tiempo de traslado \\(\tau_{hk}\\), la preparación de despegue \\(s_L\\) y la recuperación \\(s_R\\):
\\[t_k \ge t_h + \tau_{hk} + s_L \left(\sum_{l \in C,\, l \neq k} \sum_{\substack{m \in N^+ \\ \langle k,l,m \rangle \in P}} y_{klm}\right) + s_R \left(\sum_{i \in N_0,\, i \neq k} \sum_{\substack{j \in C \\ \langle i,j,k \rangle \in P}} y_{ijk}\right) - M(1 - x_{hk}) \quad \forall h \in N_0, \; k \in N^+ \setminus \{h\} \tag{16}\\]

Tiempo de llegada del dron al cliente \\(j\\) desde el nodo de despegue \\(i\\):
\\[t'_j \ge t'_i + \tau'_{ij} - M\left(1 - \sum_{\substack{k \in N^+ \\ \langle i,j,k \rangle \in P}} y_{ijk}\right) \quad \forall j \in C', \; i \in N_0 \setminus \{j\} \tag{17}\\]

Tiempo de llegada del dron al nodo de reencuentro \\(k\\), computando el tiempo de recuperación \\(s_R\\):
\\[t'_k \ge t'_j + \tau'_{jk} + s_R - M\left(1 - \sum_{\substack{i \in N_0 \\ \langle i,j,k \rangle \in P}} y_{ijk}\right) \quad \forall j \in C', \; k \in N^+ \setminus \{j\} \tag{18}\\]

El tiempo transcurrido en la salida no puede exceder la autonomía de batería \\(e\\):
\\[t'_k - (t'_j - \tau'_{ij}) \le e + M(1 - y_{ijk}) \quad \forall k \in N^+, \; j \in C \setminus \{k\}, \; i \in \{N_0 : \langle i,j,k \rangle \in P\} \tag{19}\\]

Consistencia de las variables de orden de visita \\(p_{ij}\\) del camión:
\\[u_i - u_j \ge 1 - (c+2) p_{ij} \quad \forall i, j \in C, \; i \neq j \tag{20}\\]
\\[u_i - u_j \le -1 + (c+2)(1 - p_{ij}) \quad \forall i, j \in C, \; i \neq j \tag{21}\\]
\\[p_{ij} + p_{ji} = 1 \quad \forall i, j \in C, \; i \neq j \tag{22}\\]

Impedimento de salidas solapadas (garantiza que un despegue en \\(l\\) no ocurra antes de haber recuperado el dron en \\(k\\) de la salida previa):
\\[t'_l \ge t'_k - M \left(3 - \sum_{\substack{j \in C,\, j \neq l \\ \langle i,j,k \rangle \in P}} y_{ijk} - \sum_{\substack{m \in C \setminus \{i,k,l\}}} \sum_{\substack{n \in N^+ \setminus \{i,k\} \\ \langle l,m,n \rangle \in P}} y_{lmn} - p_{il}\right) \quad \forall i \in N_0, \; k \in N^+ \setminus \{i\}, \; l \in C \setminus \{i, k\} \tag{23}\\]

### 7. Condiciones Iniciales y Naturaleza de las Variables

Condiciones de borde en el depósito:
\\[t_0 = 0 \tag{24}\\]
\\[t'_0 = 0 \tag{25}\\]
\\[p_{0j} = 1 \quad \forall j \in C \tag{26}\\]

Definición de variables binarias y no negatividad:
\\[x_{ij} \in \{0, 1\} \quad \forall i \in N_0, \; j \in N^+ \setminus \{i\} \tag{27}\\]
\\[y_{ijk} \in \{0, 1\} \quad \forall i \in N_0, \; j \in C \setminus \{i\}, \; k \in \{N^+ : \langle i,j,k \rangle \in P\} \tag{28}\\]
\\[1 \le u_i \le c+2 \quad \forall i \in N^+ \tag{29}\\]
\\[t_i \ge 0 \quad \forall i \in N \tag{30}\\]
\\[t'_i \ge 0 \quad \forall i \in N \tag{31}\\]
\\[p_{ij} \in \{0, 1\} \quad \forall i \in N_0, \; j \in C \setminus \{i\} \tag{32}\\]




\section{Conclusiones}
% Conclusiones RELEVANTES del estudio realizado. Deber\'ia responder a las preguntas: ?`todas las t\'ecnicas resuelven el mismo problema o hay algunas diferencias?, ?`En qu\'e se parecen o difieren las t\'ecnicas en el contexto del problema?, ?`qu\'e limitaciones tienen?, ?`qu\'e t\'ecnicas o estrategias son las m\'as prometedoras?, ?`existe trabajo futuro por realizar?, ?`qu\'e ideas usted propone como lineamientos para continuar con investigaciones futuras?

Se pudo observar que para el problema de TDPDP se ha realizado tanto enfoques de técnicas completas como de heuriticas y metaheuristicas. En un $53\%$ de los casos se utilizaron tecnicas heuristicas y metaheuristicas para resolver los problemas, obteniendo soluciones iniciales obtenidas con métodos completos. Esto fue debido a la naturaleza NP-Hard del problema, la cual fue evidenciada dentro del desarrollo del informe.

En base a las soluciones se evidencio la diferencia en el enfoque de como resolver el problema. Dado que el problema es un VRP modificado, se observo que una de las técnicas más utilizadas fue la de resolver primeramente un viaje para el camión y luego ir revisando en que secciones se podria cambiar la entrega por la de un dron. También resulto util diferenciar los problemas según cantidad de drones y camiones, para así enfocarse en solucionar una parte del problema que sea menos compleja que el total.

Las mayorias de limitaciones se pueden resumir a la complejidad del problema. Dado que es NP-Hard, se tarda mucho tiempo en obtener una solución del problema la cual no se vea afectada por largos tiempos de computación. Esto es debido al tamaño del espacio de búsqueda y las restricciones de sincronizar un dron con un camión. Un factor que se vio fue como las restricciones van aumentando mientras el problema se vuelve mas cercano a la realidad. En los primeros modelos, las restricciones eran muy ideales y no consideraban tanto la realidad del problema. Trabajos posteriores arreglan esto, agregando gasto de bateria variable según el peso del dron, o permitiendo que el dron se una a mitad de camino con un camión, etc.

Las tendencias muestran que el trabajo se encuentra focalizado en hacer uso de técnicas mixtas, que combinan el modelo matemático puro con las heuristicas para la resolución de los nodos utilizado spara el dron. El trabajo futuro que queda por realizar seria hacer un análisis más profundo de las diferencias entre los modelos (según cantidad de drones, camiones, batería) y como estos modelos se relacionan con la realidad de la logistica, para que así este trabajo pueda desarrollar mayor valor en empresas de entrega.


\section{Uso LLMs}
%Indique ac\'a, las herramientas de LLMs/IA usadas en su proyecto. \\En caso de no utilizar ninguna explicitar:
%\textbf{Declaro no haber utilizado ninguna herramienta de LLMs/IA en el desarrollo de este trabajo.}\\

%\noindent En caso de s\'i utilizarlo, para cada herramienta/uso, indicar:
\begin{enumerate}
\item Herramienta utilizada: Gemini Pro y NotebookLM en sus versiones gratuitas para estudiante
\item Propósito de uso: Generar las bibliografias de .bib. Revisión de redacción. Generar listados de comparación entre modelos según paper. Escritura en latex de formulas.
\item Prompts relevantes: 
\begin{itemize}
\item Generame el .bib de la referencia [X]
\item Comparame el modelo de la fuente [X] con el de la fuente [Y].
\item Dado el siguiente archivo .tex, indicame los errores ortograficos y de redacción, considerando que se debe de escribir en tercera persona.
\item Retorname en una ecuación de latex la ecuación [i] del paper [X]
\end{itemize}
\item Nivel de edición: Reescritura y revisión directa en el paper en el caso de comparaciones de modelos. Copia y pega de los .bib y formulas de latex generadas. Dependiendo del caso, las revisiones de redacción fueron seleccionadas manualmente.
\item Validación: Lectura de la referencia directa (NotebookLM entrega la sección especifica de donde saca la información).
\end{enumerate}


%Herramienta: ChatGPT (GPT-5.x)
%Uso: Generación de ideas y revisión de redacción
%Prompt principal: “Explica el algoritmo SMS-EMOA con pseudocódigo formal”
%Intervención: El texto fue adaptado, reorganizado y complementado con material de clase
%Validación: Se contrastó con [fuente X]


\pagebreak
%------Esta sección puede ser removida ----------%
\section{Bibliograf\'ia}
Indicando toda la informaci\'on necesaria de acuerdo al tipo de documento revisado. Todas las referencias deben ser citadas en el documento.
%------Remover hasta acá ----------%
\bibliographystyle{plain}
\bibliography{Referencias}

\end{document} 
